\documentclass[12pt, a4paper]{report}

% Suport pentru diacritice și alte simboluri
\usepackage{fontspec}

% Suport pentru mai multe limbi
\usepackage{polyglossia}

% Setează limba textului la română
\setdefaultlanguage{romanian}
% Am nevoie de engleză pentru rezumat
\setotherlanguages{english}

% Indentează și primul paragraf al fiecărei noi secțiuni
\SetLanguageKeys{romanian}{indentfirst=true}

% Suport pentru diferite stiluri de ghilimele
\usepackage{csquotes}

\DeclareQuoteStyle{romanian}
  {\quotedblbase}
  {\textquotedblright}
  {\guillemotleft}
  {\guillemotright}

% Utilizează biblatex pentru referințe bibliografice
\usepackage[
    maxbibnames=50,
    sorting=nty
]{biblatex}

\addbibresource{bibliography.bib}

% Setează spațiere inter-linie la 1.5
\usepackage{setspace}
\onehalfspacing

% Modificarea geometriei paginii
\usepackage{geometry}

% Include funcțiile de grafică
\usepackage{graphicx}
% Încarcă imaginile din directorul `images`
\graphicspath{{./images/}}

% Linkuri interactive în PDF
\usepackage[
    colorlinks,
    linkcolor={black},
    menucolor={black},
    citecolor={black},
    urlcolor={blue}
]{hyperref}

% Simboluri matematice codificate Unicode
\usepackage{unicode-math}

% Comenzi matematice
\usepackage{amsmath}
\usepackage{mathtools}

% Formule matematice
\newcommand{\bigO}[1]{\symcal{O}\left(#1\right)}
\DeclarePairedDelimiter\abs{\lvert}{\rvert}

% Suport pentru rezumat în două limbi
% Bazat pe https://tex.stackexchange.com/a/70818
\newenvironment{abstractpage}
  {\cleardoublepage\vspace*{\fill}\thispagestyle{empty}}
  {\vfill\cleardoublepage}
\renewenvironment{abstract}[1]
  {\bigskip\selectlanguage{#1}%
   \begin{center}\bfseries\abstractname\end{center}}
  {\par\bigskip}

% Suport pentru anexe
\usepackage{appendix}

% Stiluri diferite de headere și footere
\usepackage{fancyhdr}

\fancypagestyle{front}{
  \fancyhf{}
  \renewcommand{\headrulewidth}{0pt}
  \cfoot{}
}
\fancypagestyle{main}{
  \fancyhf{}
  \renewcommand\headrulewidth{0pt}
  \fancyhead[C]{}
  \fancyfoot[C]{\thepage}
}

% Metadate
\title{Detecția de conținut DeepFake cu ajutorul Rețelelor Neuronale Adânci}
\author{Tănase Victor-Flavian}

% Generează variabilele cu @
\makeatletter

\begin{document}

% Front matter
\cleardoublepage
\pagestyle{front}
\let\ps@plain\ps@front

% Pagina de titlu
\begin{titlepage}

% Redu marginile
\newgeometry{left=2cm,right=2cm,bottom=1cm}

\begin{figure}[!htb]
    \centering
    \begin{minipage}{0.2\textwidth}
        \includegraphics[width=\linewidth]{logo-ub.png}
    \end{minipage}
    \begin{minipage}{0.5\textwidth}
        \large
        \vspace{0.2cm}
        \begin{center}
            \textbf{UNIVERSITATEA DIN BUCUREȘTI}
        \end{center}
        \vspace{0.3cm}
        \begin{center}
            \textbf{
                FACULTATEA DE \\
                MATEMATICĂ ȘI INFORMATICĂ
            }
        \end{center}
    \end{minipage}
    \begin{minipage}{0.2\textwidth}
        \includegraphics[width=\linewidth]{logo-fmi.png}
    \end{minipage}
\end{figure}

\begin{center}
\end{center}

\vspace{1cm}

\begin{center}
\huge \textbf{\MakeUppercase{Lucrare de licență}}
\end{center}

\vspace{3cm}

\begin{center}
\large \textbf{Absolvent \\ \@author}
\end{center}

\vspace{0.25cm}

\begin{center}
\large \textbf{Coordonator științific \\ Prof. Dr. Radu Ionescu}
\end{center}

\vspace{2cm}

\begin{center}
\Large \textbf{București, iunie 2024}
\end{center}
\end{titlepage}

% Pagina 2 de titlu
\begin{titlepage}

% Redu marginile
\newgeometry{left=2cm,right=2cm,bottom=1cm}

\begin{figure}[!htb]
    \centering
    \begin{minipage}{0.2\textwidth}
        \includegraphics[width=\linewidth]{logo-ub.png}
    \end{minipage}
    \begin{minipage}{0.5\textwidth}
        \large
        \vspace{0.2cm}
        \begin{center}
            \textbf{UNIVERSITATEA DIN BUCUREȘTI}
        \end{center}
        \vspace{0.3cm}
        \begin{center}
            \textbf{
                FACULTATEA DE \\
                MATEMATICĂ ȘI INFORMATICĂ
            }
        \end{center}
    \end{minipage}
    \begin{minipage}{0.2\textwidth}
        \includegraphics[width=\linewidth]{logo-fmi.png}
    \end{minipage}
\end{figure}

\begin{center}
\end{center}

\vspace{1cm}

\begin{center}
\huge \textbf{\MakeUppercase{\@title}}
\end{center}

\vspace{3cm}

\begin{center}
\large \textbf{Absolvent \\ \@author}
\end{center}

\vspace{0.25cm}

\begin{center}
\large \textbf{Coordonator științific \\ Prof. Dr. Radu Ionescu}
\end{center}

\vspace{2cm}

\begin{center}
\Large \textbf{București, iunie 2024}
\end{center}
\end{titlepage}

\restoregeometry

\addtocounter{page}{1}

% Rezumatul in limba romana
\begin{abstractpage}

\begin{abstract}{romanian}

Evoluția capabilităților algoritmilor de inteligență artificială din ultimii ani a reinventat crearea de conținut în sfera digitală. În același timp, progresul în câmpul vederii artificiale a fost văzut de către unele entități ca o oportunitate de a răspândi dezinformare sau de a crea conținut malițios greu de detectat cu ochiul liber. 

Ceea ce a început ca cercetare academică realizată de către Justus Thies et al. \cite{thies2016face2face} a dat naștere unui mijloc de fabricare de conținut contrafăcut, capabil să manipuleze imagini și videoclipuri cu un realism nemaivăzut. Acest tip de conținut poate fi folosit in diverse scopuri cu rea-voință precum: șantaj, manipulare politică, înșelătorie, furt de identitate sau creare de conținut neadecvat, de cele mai multe ori aceste atacuri vizând persoane publice. 

În scopul combaterii dezinformării, această lucrare are ca obiectiv punerea la dispoziție către public a unui model de inteligență artificială cu o interfață web, ce poate fi ușor de utilizat pentru detectarea sau cel puțin ridicarea suspiciunii asupra imaginilor sau videoclipurilor posibil fabricate.

\end{abstract}
    
\end{abstractpage}

%Rezumatul in limba engleza
\begin{abstractpage}

\begin{abstract}{english}

The evolution of the artificial intelligence models in recent years has reinvented the creation of content on social media. In the same time, the progress in Computer Vision has been viewed by others as an opportunity to spread misinformation and create malicious content that can be hardly detected by the naked eye.

What has started as academic research by Justus Thies et al. \cite{thies2016face2face} pioneered a new way to fabricate content, capable to manipulate images and videos with unprecedented realism. This type of content can be used in many harmful ways such as blackmailing, political manipulation, fraudulent schemes, identity theft, and the creation of explicit material. Often, public figures are the primary targets of such attacks.

In order to fight disinformation, this work aims to provide the public with an artificial intelligence model with a web interface, that can be easily used to detect or at least raise suspicion over possibly fabricated images or videos. 
\end{abstract}

\end{abstractpage}

\tableofcontents

% Main matter
\cleardoublepage
\pagestyle{main}
\let\ps@plain\ps@main

\printbibliography[heading=bibintoc]

\end{document}